\documentclass[../../main.tex]{subfiles}
\begin{document}

\begin{flushright}
\footnotesize\textit{【自動文字起こし・要確認】}
\end{flushright}

{\bf [解]}
$217=7\cdot31$である．
\begin{align*}
(a-b)(a^2+ab+b^2)=7\cdot31\quad\cdots①
\end{align*}

$a^2+ab+b^2>0$から，$a-b>0$で，
\begin{align*}
(a-b,\ a^2+ab+b^2)=(1,217),\ (7,31),\ (31,7),\ (217,1)
\end{align*}

のうち，整数解を持つのは最初の2つのみで（後の2つは判別式が負），
\begin{align*}
(a,b)=(9,8),\ (-8,-9),\ (6,-1),\ (1,-6)．
\end{align*}

\newpage

\end{document}
